% TOP_PROBLEM: {"id": 30007071, "title": "Terminal numerical-dimension-one abundance in positive characteristic", "category_id": 5, "category": {"id": 5, "name": "algebraic_geometry", "display_name": "Algebraic Geometry", "description": "Geometric objects defined by polynomial equations.", "slug": "algebraic-geometry", "order_index": 5, "created_at": "2026-07-31T15:26:25.671Z"}, "status": "open"}
% FIRST_POSTED: 2026-10-06
% ATTEMPT_STATUS: solved
% OPENAI_CLAIM: {"schema_version":1,"source_repo":"https://github.com/openai/math","source_commit":"adc7f1241b42e322a6451854ab7e4b4c146bf78a","release_date":"2026-10-06","match":"full","resolution":"proved","independently_reviewed":false,"second_pass":"agreed","families":[{"family":"035","title":"Log-canonical threefold abundance in numerical dimension one","subject":"Algebraic and complex geometry","manuscripts":[{"title":"Abundance in numerical dimension one for terminal threefolds in positive characteristic","pdf_path":"preprints/Abundance-in-numerical-dimension-one-for-terminal-threefolds-in-positive-characteristic-September-24-2026/paper.pdf","readme_path":"preprints/Abundance-in-numerical-dimension-one-for-terminal-threefolds-in-positive-characteristic-September-24-2026/README.md","date":"2026-09-24","theorem_ref":"thm:main (sections/introduction.tex:14)","statement_sources":[{"path":"preprints/Abundance-in-numerical-dimension-one-for-terminal-threefolds-in-positive-characteristic-September-24-2026/build/sections/introduction.tex","line":14,"label":"thm:main"}]}],"lean":null,"reasoning_trace":null}]}
\documentclass[11pt]{article}
\usepackage[margin=25mm]{geometry}
\usepackage{fontspec}
\setmainfont{Latin Modern Roman}
\usepackage{amsmath,amssymb,mathtools}
\usepackage{unicode-math}
\setmathfont{Latin Modern Math}
\usepackage{enumitem}
\usepackage{xurl,hyperref,bookmark}
\hypersetup{colorlinks=true,urlcolor=blue}
\setcounter{secnumdepth}{0}
\setlength{\parindent}{0pt}
\setlength{\parskip}{5pt}
\setlength{\emergencystretch}{3em}
\begin{document}
\section*{Terminal numerical-dimension-one abundance in positive characteristic}

\subsection{OpenAI's claim}
For terminal factorial projective threefolds in characteristic $p>3$, a nef canonical divisor of numerical dimension one is claimed to have Kodaira dimension one and be semiample. The cited manuscript argues for that result; its relation to this problem is specified below.

\subsection{Mathematical statement of the claimed result}
\paragraph{Mathematical claim 1.} Over an algebraically closed field of characteristic $p>3$, let $X$ be a normal projective rational-factorial terminal threefold with nef $K_X$. If an ample Cartier $H$ satisfies $K_XH^2>0$ and $K_X^2H=0$, OpenAI claims
\[
\exists m>0:\ mK_X\text{ Cartier},\ h^0(X,mK_X)\ge2,\qquad\kappa(X,K_X)=1,
\]
and $K_X$ semiample. This is the boundary-zero terminal numerical-dimension-one theorem in the stated prime range.

\subsection{Outline of the argument}
Manuscript 1 (family 035, 2026-09-24): The proof starts from known positive-characteristic nonvanishing and seeks a second independent canonical section. A primitive-cycle obstruction combines tangent-sheaf foliations, strongly semistable inequalities, surface abundance, and Cartier descent; the resulting positive Kodaira dimension invokes the existing semiampleness theorem.

\subsection{Scope relative to this problem}
Manuscript 1 (family 035, 2026-09-24): The problem selects the mathematical statement summarized here, including its stated hypotheses and parameter range. The cited theorem gives a claimed proof of that assertion in full.

\subsection{Formal verification}
Family 035 has no Lean scope page in this snapshot; no formalization of these manuscript results is listed.

This site has not executed the Lean files or checked their dependencies. These descriptions report the scope of the linked declarations, rather than a new verification of the proofs.

\subsection{Status caveat}
OpenAI attributes these results to an internal model. Its release notice cautions that results lacking full formalization can contain mistakes. This site has not independently assessed the proofs; a solved label records the scope of the reported claim.

\subsection{Sources}
\begin{itemize}
\item OpenAI, \emph{Abundance in numerical dimension one for terminal threefolds in positive characteristic}, 2026-09-24; citation key \texttt{OAI:Abundance-in-numerical-dimension-one-for-terminal-threefolds-in-positive-characteristic-September-24-2026}. \href{https://github.com/openai/math/blob/adc7f1241b42e322a6451854ab7e4b4c146bf78a/preprints/Abundance-in-numerical-dimension-one-for-terminal-threefolds-in-positive-characteristic-September-24-2026/paper.pdf}{PDF (GitHub)}; \href{https://github.com/openai/math/raw/adc7f1241b42e322a6451854ab7e4b4c146bf78a/preprints/Abundance-in-numerical-dimension-one-for-terminal-threefolds-in-positive-characteristic-September-24-2026/paper.pdf}{PDF (direct)}; \href{https://github.com/openai/math/blob/adc7f1241b42e322a6451854ab7e4b4c146bf78a/preprints/Abundance-in-numerical-dimension-one-for-terminal-threefolds-in-positive-characteristic-September-24-2026/README.md}{Manuscript page and citation}; \href{https://github.com/openai/math/blob/main/preprints/Abundance-in-numerical-dimension-one-for-terminal-threefolds-in-positive-characteristic-September-24-2026/paper.pdf}{Latest version}; Statement source: \href{https://github.com/openai/math/blob/adc7f1241b42e322a6451854ab7e4b4c146bf78a/preprints/Abundance-in-numerical-dimension-one-for-terminal-threefolds-in-positive-characteristic-September-24-2026/build/sections/introduction.tex\#L14}{thm:main}.
\item \href{https://github.com/openai/math/blob/adc7f1241b42e322a6451854ab7e4b4c146bf78a/CONTENTS.md}{OpenAI catalogue entry}.
\end{itemize}
\end{document}
